\bigskip

\textbf{\large Partie I}

\medskip

Dans un repère orthogonal d'origine O, on considère les points A(0 ; 12), B(1 ; 12) et C(1 ; 0). Dans la figure ci-dessous, deux courbes $\mathcal{C}_1$ et $\mathcal{C}_2$ délimitent un domaine grisé $\mathcal{D}_1$, un domaine hachuré $\mathcal{D}_2$, et un domaine D$\mathcal{ }_3$ non coloré, à l'intérieur du rectangle OABC.

\begin{center}
\psset{xunit=9.2cm,yunit=0.4583,comma=true,labelFontSize=\scriptstyle}
\begin{pspicture}(-0.1,-1)(1.3,14)
\psline[linewidth=0.55pt](0,12)(1,12)\psline[linewidth=0.55pt](1,0)(1,12)

 \psaxes[linewidth=0.55pt,Dx=0.25,Dy=3]{-}(1.1,12.2)

 \def\Func{x 2 exp x 1 add mul 6 mul}
   \psplot[plotpoints=1000,linewidth=0.95pt,linecolor=blue]{0}{1}{\Func}
 \pscustom[fillstyle=solid,fillcolor=lightgray]{\psplot{0}{1}{x 2  exp x 1 add mul 6 mul}, \psline(1,12)(1,0)}

\pscustom[fillstyle=hlines,hatchwidth=0.55pt,hatchcolor=red, linecolor=red]{\pscurve(0,3)(0.25,4.5)(0.5,7.5)(0.75,12)\psline(0.75,12)(0,12)}
\psline(1,12)(0,12)
\psline(1,0)(0,0)
\pscurve[linewidth=0.45pt,linecolor=red](0,3)(0.25,4.5)(0.5,7.5)(0.75,12)

\uput[d](0.3,9){$\mathcal{D}_2$} \uput[d](0.55,5.5){$\mathcal{D}_3$} \uput[u](0.79,2.5){$\mathcal{D}_1$}
\uput[r](0.01,2.6){$\mathcal{C}_1$}\uput[r](0.88,11){$\mathcal{C}_2$}
\uput[u](0.02,12){\blue A}\uput[ur](1,12){\blue B}\uput[ur](1,0){\blue C}\uput[dl](0,0){\blue O}
\end{pspicture}
\end{center}

La courbe $\mathcal{C}_2$ est la représentation graphique de la fonction $f $définie sur l'intervalle [0;1] par:

\[f(x) = 6x^3 + 6x^2\]

\begin{enumerate}
\item  Calculer l'intégrale $\int_0^1f(x)\mathrm{d}x$
\end{enumerate}

On admet que la fonction $f$ est positive sur l'intervalle [0 ; 1]. 
\begin{enumerate}[resume]
\item Interpréter le résultat du calcul précédent dans le contexte de l'exercice.
\end{enumerate}

L'aire du domaine $\mathcal{D}_2$ hachuré a été calculée (en unité d'aire) à l'aide d'un logiciel et vaut $\dfrac{135}{32}$.
\begin{enumerate}[resume]
\item  Estimer l'aire en unité d'aire du domaine $\mathcal{D}_3$.
\end{enumerate}

\bigskip

\textbf{\large Partie II}

\medskip

On considère les nombres complexes $z_1 = 1-\i\sqrt{3}$ et $z_2= 4\e^{\i\frac{\pi}{6}}$ où $\i$ désigne le nombre complexe de module 1 et d'argument $\dfrac{\pi}{2}$.

\begin{enumerate}
\item  Écrire le nombre $z_1$ sous forme exponentielle. Détailler les calculs.

\item Démontrer que le nombre $z = z_1^2 \times  z_2$, est un nombre imaginaire pur en détaillant les calculs.
\end{enumerate}

\bigskip
